\documentclass{article}
\usepackage{geometry}
\usepackage{booktabs}
\usepackage{hyperref}
\geometry{margin=1in}
\begin{document}
\title{Genetic Programming Classifier Report}
\author{COS314 Project}
\date{\today}
\maketitle

\begin{abstract}
This report summarizes the implementation and comparison of two Genetic Programming classifiers on the Breast 
Cancer Wisconsin dataset. One variant evolves arithmetic expressions, while the other evolves tree-based decision
\end{abstract}

\section{Introduction}
Genetic Programming (GP) evolves classifiers by applying selection, crossover, and mutation to a 
population of program trees. 
This project implements two GP variants in Java and evaluates them on a 
standard breast cancer dataset.

\subsection{Experimental Setup}
Each algorithm was executed for 30 independent runs using seeds given.
The best-performing run was selected based on training fitness and saved. Saved in results directory. 
The seed used for the demonstrated Decision Tree GP run was \texttt{0}, also the case for the demonstrated Arithmetic GP.

\section{Methods}
The two variants are:
\begin{itemize}
  \item arithmetic-expression GP, which evolves symbolic numerical formulas
  \item decision-tree GP, which evolves logical tree-based classifiers
\end{itemize}
Both variants use the same dataset and evaluation metrics.
Each variant is ran 30 times with the same inputted seeds.
The program saves one result file per run under \texttt{results/}.

\subsection{GP Parameters}
\begin{center}
    \begin{tabular}{lc}
        \toprule
        Parameter & Value \\
        \midrule
        Population Size & 200 \\
        Maximum Generations & 100 \\
        Crossover Rate & 0.8 \\
        Mutation Rate & 0.2 \\
        Maximum Depth & 5 \\
        Tournament Size & 5 \\
        Selection Method & Tournament Selection \\
        Initialisation Method & Ramped Half-and-Half \\
        Fitness Function & Training Accuracy \\
        \bottomrule
    \end{tabular}
\end{center}

\section{Results}
The main comparison is based on accuracy, F-measure, and runtime.
In the tested runs, the decision-tree GP achieved higher test F-measure, while the arithmetic GP reached a bit higher training accuracy.

\begin{center}
\begin{tabular}{lcccc}
\toprule
Algorithm & Training (\%) & Test (\%) & F-measure & Runtime (s) \\
\midrule
Decision Tree GP & 95.08 & 53.49 & 0.3750 & 0.49 \\
Arithmetic GP & 95.08 & 50.00 & 0.3944 & 0.35 \\
\bottomrule
\end{tabular}
\end{center}

A two sample t-test on test accuracies found no statistically significant difference at 
the 0.05 level for this experimental setup.

\subsection{Program Output}
During training, the program prints the best individual of each generation together with its training accuracy.

\section{Discussion}
The decision-tree variant appears to generalize better on the test split, while the arithmetic variant 
fits the training set more closely.
Further work should include cross-validation, parameter tuning and comparison with standard machine learning strategies.

\section{Although}
limitation observed during testing was occasional overfitting, where individuals achieved 
high training accuracy but weaker test performance. 
This was more noticeable in the arithmetic GP variant.

\section{Conclusion}
This project implements two GP classifiers and compares their performance on breast cancer classification.

\end{document}
